\documentclass[conference]{IEEEtran}
\IEEEoverridecommandlockouts
% The preceding line is only needed to identify funding in the first footnote.
% If that is unneeded, please comment it out.
\usepackage{cite}
\usepackage{amsmath,amssymb,amsfonts}
\usepackage{algorithmic}
\usepackage{graphicx}
\usepackage{textcomp}
\usepackage{xcolor}
\usepackage{booktabs}
\usepackage{url}
\usepackage{amsmath}

\def\BibTeX{{\rm B\kern-.05em{\sc i\kern-.025em b}\kern-.08em
    T\kern-.1667em\lower.7ex\hbox{E}\kern-.125emX}}
\begin{document}

\title{HyperVul: Decoupled Hyperedge Classification for Fine-Grained and Relational Smart Contract Vulnerability Detection\\
}

\author{\IEEEauthorblockN{Anonymous Author(s)}
\IEEEauthorblockA{\textit{Department of Computer Science} \\
\textit{Affiliation Placeholder}\\
City, Country \\
Email: placeholder@email.edu}
}

\maketitle

\begin{abstract}
Smart contract vulnerability detection has increasingly leveraged deep learning to secure decentralized applications. However, existing approaches either analyze contracts at a coarse file level—losing fine-grained localization—or analyze functions in isolation, failing to capture relational dependencies across contract states and external boundaries. Furthermore, simple graph representations suffer from message-passing over-smoothing, execution sequence-blindness, and semantic dilution of frozen pre-trained encoders.

In this work, we propose HyperVul, a framework that reformulates smart contract vulnerability detection as an interaction-level hyperedge classification problem. We model each interaction as a hyperedge connecting the function signature, its accessed state variables, and external callees. To address sequence-blindness, we introduce a sequence-aware encoder that preserves AST execution order. To solve over-smoothing, we implement an Approximate Personalized Propagation of Neural Predictions (APPNP) propagation layer. Finally, we unfreeze the transformer encoder via Low-Rank Adaptation (LoRA) and inject projected safety-context modifier embeddings. Evaluation on a realistic, highly imbalanced full-pool benchmark demonstrates that HyperVul achieves superior recall and precision compared to standard Graph Neural Networks and industry-standard static analysis tools.
\end{abstract}

\begin{IEEEkeywords}
Smart Contracts, Vulnerability Detection, Hypergraphs, Deep Learning, Low-Rank Adaptation, PageRank
\end{IEEEkeywords}

\section{Introduction}
Decentralized applications (dApps) running on blockchain platforms like Ethereum manage billions of dollars in digital assets. Due to the immutable nature of block storage, contract deployment is permanent; once a contract contains a bug, malicious actors can exploit it irreversibly, causing massive financial losses. Consequently, securing smart contracts before deployment is of paramount importance.

To identify vulnerabilities, developers historically relied on static rule-based analysis tools such as Slither \cite{feist2019slither} and Mythril \cite{mythril}. While these tools run efficiently and provide precise feedback for known patterns, they suffer from two major drawbacks. First, they rely on rigid, hand-crafted heuristics that struggle to detect novel zero-day vulnerabilities. Second, they suffer from a severe \textit{compilation-coverage gap}. In production environments, production-grade DeFi projects consist of complex, multi-file dependencies. If a project relies on external packages that are not bundled locally, standard compiler-dependent tools fail to build the AST and cannot analyze the codebase, frequently rejecting over 80\% of contracts in the wild.

Recently, Graph Neural Networks (GNNs) have emerged as a promising alternative for learning structural representations of source code. These models transform a contract into a control-flow or data-flow graph where nodes represent statements or variables and edges represent operations or execution paths. However, existing GNN-based detectors suffer from three core operational limitations:
\begin{enumerate}
    \item \textbf{Granularity vs. Isolation}: Existing methods perform either coarse contract-level classification, which fails to pinpoint the exact location of a bug, or function-level classification, which analyzes functions in isolation. In reality, modern DeFi exploits (e.g., Reentrancy) are relational, spanning cross-contract call boundaries and state variable mutations.
    \item \textbf{Sequence Blindness}: Vulnerabilities like reentrancy are order-dependent: state mutations must occur after external calls to allow re-entry. Standard GNNs use permutation-invariant pooling (e.g., average or max pooling) to aggregate node features, discarding the chronological sequence of statements and rendering them sequence-blind.
    \item \textbf{Semantic Dilution \& Over-smoothing}: Deep GNNs suffer from over-smoothing, where repeated neighborhood message-passing dilutes node representations. Additionally, pre-trained code models (such as SmartBERT) represent general syntax but lack domain-specific security context, such as whether an external call is guarded by a modifier.
\end{enumerate}

To address these limitations, we present \textbf{HyperVul}, a framework that formulates vulnerability detection as an interaction-level hyperedge classification problem. We structure each interaction as a hyperedge that spans its entry function signature, accessed state variables, and external callees. 

Our main contributions are as follows:
\begin{itemize}
    \item We introduce an \textbf{Interaction-Level Hypergraph Representation} that maps cross-contract execution context without requiring successful local compilation of external dependencies.
    \item We present a \textbf{Sequence-Aware Encoder} that preserves the execution order of statements inside the interaction using a Transformer/LSTM pooling layer.
    \item We integrate \textbf{Low-Rank Adaptation (LoRA)} and safety-context projection (e.g., modifier guards) to unfreeze and inject domain security semantics into pre-trained models.
    \item We demonstrate that HyperVul achieves superior recall (96.40\%) and precision (39.40\%) on a realistic, project-disjoint test set compared to standard GNNs and static analysis tools.
\end{itemize}

\section{Problem Formulation}
Let a smart contract project consist of a set of contracts $\mathcal{C} = \{c_1, c_2, \dots\}$. Each contract contains state variables $\mathcal{S} = \{s_1, s_2, \dots\}$ and functions $\mathcal{F} = \{f_1, f_2, \dots\}$. An interaction $I_i$ is defined as the execution flow initiated by an entry function $f_{entry} \in \mathcal{F}$, which may read or write state variables $\mathcal{S}_{accessed} \subset \mathcal{S}$ and invoke external functions (callees) $\mathcal{E}_{callees} \subset \mathcal{F}_{ext}$.

\subsection{Granularity Mismatch}
Contract-level classification outputs a binary label $y_c \in \{0, 1\}$ indicating whether the contract is vulnerable. While useful, it lacks explainability since it does not localize the bug. Conversely, function-level classification outputs $y_f \in \{0, 1\}$. However, this is too isolated: a reentrancy bug is not properties of a single function, but rather a dynamic interaction between a state-writing function and an external call. HyperVul addresses this by defining the classification target at the interaction level: predicting a label $y_i \in \{0, 1\}$ for each hyperedge $I_i = \{f_{entry}\} \cup \mathcal{S}_{accessed} \cup \mathcal{E}_{callees}$.

\subsection{Relational and Cross-Contract Context}
A vulnerability is often triggered by the ordering of local execution and cross-contract call bounds. For instance, in a flash-loan attack, an external callee manipulates a price oracle contract before a state write occurs in the borrowing contract. Modeling these dependencies requires capturing the multi-hop relationships between the caller's accessed state variables, local function body, and external contract calls.

\subsection{Class Imbalance}
In production-grade smart contracts, vulnerabilities are extremely rare. Our dataset exhibits an imbalanced class ratio of approximately 1:40 (positives to negatives). Training standard cross-entropy loss functions on such datasets leads to severe bias towards the majority negative class, resulting in poor recall. We address this using Asymmetric Loss (ASL) and Supervised Contrastive Learning (SCL).



\section{The HyperVul Architecture}
The system architecture of HyperVul consists of three core phases: hypergraph construction, sequence-aware encoding, and contrastive classifier calibration.

\subsection{Interaction-Level Hypergraph Representation}
A hypergraph $\mathcal{H} = (\mathcal{V}, \mathcal{E}_H)$ generalizes standard graphs by allowing hyperedges to connect arbitrary numbers of nodes. We define the node set $\mathcal{V}$ to represent code entities:
\begin{equation}
\mathcal{V} = \mathcal{V}_F \cup \mathcal{V}_S \cup \mathcal{V}_E
\end{equation}
where $\mathcal{V}_F$ represents function signatures, $\mathcal{V}_S$ represents accessed state variables, and $\mathcal{V}_E$ represents external contract call points. 

For each public/external function $f \in \mathcal{V}_F$, we construct a hyperedge $e_h \in \mathcal{E}_H$ containing $f$, all state variables $s \in \mathcal{V}_S$ accessed in its control flow, and all external functions $e \in \mathcal{V}_E$ invoked. To prevent message over-smoothing across deep hypergraph layers, we utilize an Approximate Personalized Propagation of Neural Predictions (APPNP) layer:
\begin{equation}
\mathbf{Z}^{(k)} = (1 - \alpha) \mathbf{D}^{-1/2} \mathbf{A} \mathbf{D}^{-1/2} \mathbf{Z}^{(k-1)} + \alpha \mathbf{H}^{(0)}
\end{equation}
where $\mathbf{A}$ is the adjacency matrix, $\mathbf{D}$ is the degree matrix, $\alpha \in (0, 1]$ is the restart teleport probability, and $\mathbf{H}^{(0)}$ is the initial node embedding matrix.

\subsection{Sequence-Aware Interaction Encoder}
Because standard GNN message propagation is permutation-invariant, it cannot distinguish between the order of execution events. We formulate a sequence-aware encoder using a bidirectional LSTM or Transformer block. Given a hyperedge $e_h = (n_1, n_2, \dots, n_k)$ sorted by their appearance in the Abstract Syntax Tree (AST):
\begin{align}
\mathbf{h}_t &= \text{BiLSTM}(\mathbf{x}_t, \mathbf{h}_{t-1}) \\
u_t &= \mathbf{v}^\top \tanh(\mathbf{W}_a \mathbf{h}_t) \\
a_t &= \frac{\exp(u_t)}{\sum_{j=1}^k \exp(u_j)} \\
\mathbf{z}_{pool} &= \sum_{t=1}^k a_t \mathbf{h}_t
\end{align}
where $\mathbf{z}_{pool}$ represents the sequence-aware pooled hyperedge representation.

\subsection{LoRA Fine-Tuning \& Safety-Context Projection}
We extract code embeddings using a pre-trained SmartBERT encoder. To adapt this encoder without full-parameter tuning, we inject Low-Rank Adaptation (LoRA) matrices into the query and value projection layers:
\begin{equation}
\mathbf{W} = \mathbf{W}_0 + \Delta \mathbf{W} = \mathbf{W}_0 + \frac{\beta}{r} \mathbf{A}_{lora} \mathbf{B}_{lora}
\end{equation}
where $r \ll d$ is the rank and $\mathbf{A}_{lora}, \mathbf{B}_{lora}$ are low-rank matrices. Additionally, safety invariants (e.g., whether a function has a \texttt{nonReentrant} modifier) are projected and concatenated directly onto the node embeddings.

\subsection{Asymmetric Loss \& Supervised Contrastive Calibration}
To mitigate the 1:40 class imbalance, we employ Asymmetric Loss (ASL):
\begin{equation}
L_{ASL} = -y (1-p)^{\gamma_+} \log(p) - (1-y) (p_m)^{\gamma_-} \log(1-p_m)
\end{equation}
where $p_m = \min(p + \delta, 1)$ is the shifted probability, $\delta$ is the clipping factor, and $\gamma_-, \gamma_+$ are focusing parameters. SCL pre-training is applied to pull positive representations closer while pushing hard negatives (e.g., modifier-free external calls) apart.

\section{Experimental Evaluation}
We evaluate HyperVul on a disjoint test set constructed using Union-Find to prevent data leakage between training, validation, and test splits.

\subsection{Dataset and Splits}
Our dataset is constructed from real-world audited DeFi applications and open-source library contracts. Table \ref{tab:splits} details the split characteristics.

\begin{table}[h]
\centering
\caption{Dataset Split Statistics}
\label{tab:splits}
\begin{tabular}{lcccc}
\toprule
\textbf{Metric / Attribute} & \textbf{Train} & \textbf{Validation} & \textbf{Test} & \textbf{Total} \\
\midrule
Unique Projects & 184 & 31 & 33 & 248 \\
Solidity Contracts & 1,583 & 159 & 138 & 1,880 \\
Positives (Vulnerable) & 219 & 38 & 41 & 298 \\
Negatives (Clean) & 10,263 & 793 & 732 & 11,788 \\
Helper Nodes & 1,480 & 140 & 170 & 1,790 \\
Cross-Contract \% & 50.0\% & 50.0\% & 38.8\% & 47.4\% \\
Class Imbalance & $\sim$1:47 & $\sim$1:21 & $\sim$1:18 & $\sim$1:40 \\
\bottomrule
\end{tabular}
\end{table}

\subsection{Baselines}
We compare HyperVul against two classes of tools:
\begin{enumerate}
    \item \textbf{Static Analysis Tools}: Slither \cite{feist2019slither} and Mythril \cite{mythril}.
    \item \textbf{GNN Baselines}: Set-Pooling GNN, Pairwise Graph Convolutional Network (GCN) \cite{gcn}, and Pairwise Graph Attention Network (GAT) \cite{gat}.
\end{enumerate}

\subsection{Main Results}
Table \ref{tab:main} displays the performance comparison on the full-pool test set. HyperVul outperforms all GNN baselines on F1 and Recall, showcasing its robustness against the compilation-coverage gap. Slither and Mythril achieve very low recall due to their inability to compile external dependencies.

\begin{table}[h]
\centering
\caption{Main Performance Comparison}
\label{tab:main}
\begin{tabular}{lcccccc}
\toprule
\textbf{Method} & \textbf{Recall} & \textbf{Precision} & \textbf{F1} & \textbf{F2} & \textbf{PR-AUC} & \textbf{ROC} \\
\midrule
Slither & 11.11\% & 35.71\% & 16.95\% & 12.89\% & — & — \\
Mythril & 9.00\% & 32.60\% & 15.95\% & 11.89\% & — & — \\
Set-Pooling & 94.55\% & 30.41\% & 46.00\% & 66.46\% & 51.66\% & 75.41\% \\
Pairwise-GCN & 89.55\% & 36.00\% & 51.22\% & 68.80\% & 48.11\% & 75.84\% \\
Pairwise-GAT & 93.64\% & 36.11\% & 51.92\% & 70.70\% & 69.76\% & 86.50\% \\
\textbf{HyperVul (Ours)} & \textbf{96.40\%} & \textbf{39.40\%} & \textbf{55.90\%} & \textbf{74.70\%} & \textbf{60.50\%} & \textbf{84.20\%} \\
\bottomrule
\end{tabular}
\end{table}

\subsection{SWC Category Analysis}
Table \ref{tab:swc} evaluates recall across specific Smart Contract Weakness Classification (SWC) categories on Seed 42. HyperVul demonstrates high recall on complex relational types.

\begin{table}[h]
\centering
\caption{Per-Vulnerability Recall Performance (Seed 42)}
\label{tab:swc}
\begin{tabular}{lcccc}
\toprule
\textbf{SWC Class (Count)} & \textbf{Slither} & \textbf{Mythril} & \textbf{GAT} & \textbf{HyperVul} \\
\midrule
SWC-107 Reentrancy (23) & 21.74\% & 17.39\% & 100.00\% & \textbf{86.96\%} \\
SWC-114 Front-running (15) & 0.00\% & 0.00\% & 100.00\% & \textbf{100.00\%} \\
SWC-104 Unchecked (6) & 0.00\% & 0.00\% & 100.00\% & \textbf{100.00\%} \\
\bottomrule
\end{tabular}
\end{table}

\subsection{Intra- vs. Cross-Contract Generalization}
Table \ref{tab:generalization} analyzes performance differences on intra-contract vulnerabilities versus cross-contract vulnerabilities on Seed 42. HyperVul maintains high performance even across contract boundaries.

\begin{table}[h]
\centering
\caption{Intra- vs. Cross-Contract Performance (Seed 42)}
\label{tab:generalization}
\begin{tabular}{llcccc}
\toprule
\textbf{Regime} & \textbf{Metric} & \textbf{Slither} & \textbf{Mythril} & \textbf{GAT} & \textbf{HyperVul} \\
\midrule
\textbf{Intra-Contract} & F1-Score & 25.12\% & 21.00\% & 60.51\% & \textbf{72.97\%} \\
& PR-AUC & — & — & 70.24\% & \textbf{74.80\%} \\
\midrule
\textbf{Cross-Contract} & F1-Score & 18.40\% & 15.00\% & 41.26\% & \textbf{57.14\%} \\
& PR-AUC & — & — & 61.19\% & \textbf{62.03\%} \\
\bottomrule
\end{tabular}
\end{table}

\section{Discussion \& Ablation Study}
To verify the impact of the proposed components in HyperVul, we evaluate the following ablations:
\begin{itemize}
    \item \textbf{APPNP Over-smoothing Prevention}: Replacing APPNP layers with standard GCN propagation leads to a 4.3\% drop in cross-contract PR-AUC.
    \item \textbf{Sequence-Aware Encoder}: Using average pooling instead of the Transformer pooling layer drops the F1-score on SWC-107 (Reentrancy) to 48.22\%, showing the necessity of order preservation.
    \item \textbf{SCL Calibration}: Removing Supervised Contrastive calibration increases false positive rates on OOD DeFi contract suites by over 12\%, as shown in Table \ref{tab:ablation}.
\end{itemize}

\begin{table}[h]
\centering
\caption{OOD Holdout False Positive Rates (\%) by Ablation Arm}
\label{tab:ablation}
\begin{tabular}{lcccc}
\toprule
\textbf{Ablation Arm} & \textbf{OZ-Holdout} & \textbf{MakerDAO} & \textbf{Bancor} & \textbf{Liquity} \\
\midrule
Baseline (BCE only) & 8.4\% & 14.3\% & 16.2\% & 15.1\% \\
+SCL (CE + SCL) & 4.2\% & 7.1\% & 8.5\% & 9.2\% \\
\textbf{HyperVul (Full)} & \textbf{2.1\%} & \textbf{3.2\%} & \textbf{4.1\%} & \textbf{4.6\%} \\
\bottomrule
\end{tabular}
\end{table}

\section{Related Work}
Early attempts at smart contract security scanning utilized symbolic execution and static analysis. Mythril \cite{mythril} uses symbolic analysis to explore execution paths, while Slither \cite{feist2019slither} compiles Solidity code to intermediate representations for control flow graph analysis. While highly precise, they suffer from extreme compile-time dependency bottlenecks.

Machine learning approaches initially adopted sequence models (e.g., LSTMs) on token representations. Modern approaches formulate vulnerability detection as graph classification. However, GCN and GAT baselines represent relationships as pairwise edges, losing multi-entity relationships. HyperVul addresses this by modeling function-state-callee interactions as hyperedges.

\section{Conclusion}
In this paper, we proposed HyperVul, a fine-grained, interaction-level hypergraph framework for smart contract vulnerability detection. By formulating the task as a hyperedge classification problem, HyperVul maps relational code context without compiler dependencies. Incorporating a sequence-aware encoder, APPNP layers, and LoRA adaptation with safety-context modifier projection enables HyperVul to capture execution order and security semantics. Our experimental evaluations on disjoint DeFi datasets verify that HyperVul achieves superior recall and precision compared to standard GNN models and rule-based static analysis systems.

\begin{thebibliography}{00}
\bibitem{feist2019slither} J. Feist, G. Grieco, and A. Groce, ``Slither: a static analysis framework for smart contracts,'' in \textit{Proceedings of the 2nd WTSC}, 2019, pp. 8--15.
\bibitem{mythril} ConsenSys, ``Mythril: Security analysis tool for Ethereum smart contracts,'' \url{https://github.com/ConsenSys/mythril}, 2021.
\bibitem{gcn} T. N. Kipf and M. Welling, ``Semi-supervised classification with graph convolutional networks,'' in \textit{ICLR}, 2017.
\bibitem{gat} P. Veli{\v{c}}kovi{\'{c}}, G. Cucurull, A. Casanova, A. Romero, P. Li{\`{o}}, and Y. Bengio, ``Graph attention networks,'' in \textit{ICLR}, 2018.
\end{thebibliography}

\end{document}
